\section{Motivating Example}
\label{sec:motivating-example}

To concretize the challenge of semantically blind heuristics and to illustrate why the expert-like intuition introduced in the introduction is effective, we present a motivating example from a real-world program analysis scenario.

\subsection{A Challenging SMT Formula}

Consider the SMT formula, a conjunction of four constraints, defined over six integer variables: $x, y, z, a, b, c$. This formula is representative of path conditions generated by symbolic execution of programs with complex data transformations. The full formula is shown in Figure~\ref{fig:motivating-smt}.
\wyc{The structure of this paragraph is incorrect. First say "Figure~\ref{fig:motivating-smt} presents a real-world SMT formula, which is simplified from a path constraint collected from xxx program.". Then, explain the formula, \eg, "it consists of six integer variable...". The next paragraph only discusses why this formula is challenging.}


\begin{figure}[t]
\centering
\Description{A framed SMT-LIB style listing illustrating nonlinear and modular constraints (C1--C4).}
\caption{An illustrative SMT formula with nonlinear integer constraints.}
\label{fig:motivating-smt}
\begin{lstlisting}[language={},
  basicstyle=\ttfamily\normalsize,
  frame=single,
  framerule=0.4pt,
  framesep=4pt,
  framexrightmargin=10pt,
  xrightmargin=10pt,
  breaklines=false,
  breakatwhitespace=false,
  aboveskip=0.5pt,
  belowskip=0.5pt,
  columns=fullflexible,
  keepspaces=true,
  captionpos=b]
C1: (((x * 37 + 23) % 101) + ((x - 30)^2 % 57) + y^3 - z^2 + a^2 - b + c) == 150
C2: (x^4 + y^3 - 5*z*a + b - c^2) < 400
C3: ((x^2 - y^2) + ((z*12 + 301) % 101) + ((z-15)*(z-30) % 57) + 
    ((z+15)*(z-30)*(z-11) % 9) + a - b*c) > 
    (((a*12 + 301) % 101) + ((a-15)*(a-30) % 57) + ((a+15)*(a-30)*(a-11) % 9))
C4: (((x^2 - y^2) + ((z*37 + 23) % 101) + ((z - 30)^2 % 57) + a - b*c)) != 55
// Full constraint: C1 AND C2 AND C3 AND C4
\end{lstlisting}
\end{figure}

This formula exhibits several sources of complexity that are notoriously difficult for SMT solvers. High-degree nonlinear polynomials (\eg, $x^4$) and nested products such as $(z+15)(z-30)(z-11)$ inflate a vast, nonconvex search space, consistent with the inherent hardness of nonlinear integer programming~\cite{murty2002complexity}. Modular arithmetic (the operator \texttt{\%}) introduces discontinuities that force solvers into large case splits, often yielding exponential blow-up. Moreover, variables are strongly coupled across multiple, interdependent constraints; a change to one variable propagates throughout the formula, undermining local propagation and pruning. As a result, even a state-of-the-art solver like Z3~\cite{demoura2008z3} fails to return a result within a 1200-second timeout on the same workstation as used in Section~\ref{sec:evaluation}, exemplifying the solver bottleneck that motivates our approach. \wyc{As a result, even a state-of-the-art solver like Z3~\cite{demoura2008z3} fails to return a result within a 1200-second timeout according to our study (Section~\ref{sec:evaluation})....}

% \subsection{Strategic Variable Analysis}
% To systematically identify which variable should be concretized, we can assess each variable's potential impact by analyzing its structural role in the formula. We propose a heuristic framework based on three metrics:
% \begin{itemize}
%     \item \textbf{Frequency (F)}: The total number of times a variable appears.
%     \item \textbf{Complexity (C)}: The number of times a variable is involved in high-cost nonlinear operations (e.g., powers, products of variables, modulo).
%     \item \textbf{Coupling (P)}: The number of distinct constraints (C1-C4) in which the variable participates.
% \end{itemize}

% We can score each variable based on these metrics, with a higher total "Impact Score" ($I = F+C+P$) indicating a greater potential for simplification upon concretization. Table~\ref{tab:impact-score} shows this analysis for our example.

% \begin{table}[h]
% \centering
% \caption{Heuristic impact scores for each variable in the example.}
% \label{tab:impact-score}
% \begin{tabular}{@{}lcccc@{}}
% \toprule
% \textbf{Variable} & \textbf{Frequency (F)} & \textbf{Complexity (C)} & \textbf{Coupling (P)} & \textbf{Impact Score (I)} \\ \midrule
% $x$ & 3 & 2 ($x^4$, $x^2$) & 3 (C1,C2,C4) & 8 \\
% $y$ & 3 & 2 ($y^3$, $y^2$) & 3 (C1,C2,C3) & 8 \\
% $\mathbf{z}$ & \textbf{4} & \textbf{3 ($z^2$, mod, nested)} & \textbf{4 (C1,C2,C3,C4)} & \textbf{11} \\
% $a$ & 4 & 2 ($a^2$, mod) & 3 (C1,C2,C3) & 9 \\
% $b$ & 2 & 1 ($b*c$) & 3 (C1,C2,C3) & 6 \\
% $c$ & 2 & 2 ($c^2$, $b*c$) & 2 (C1,C2) & 6 \\ \bottomrule
% \end{tabular}
% \end{table}

% The analysis clearly identifies variable $\mathbf{z}$ as the most influential, with the highest Impact Score of 11. It is the only variable present in all four constraints and is involved in the most complex operations. This makes it the prime candidate for concretization.

\subsection{Guided Simplification via Variable Concretization}

Our strategy is to simplify the formula by first identifying and then concretizing the most influential variable. To select the target\wyc{what is target?}, we employ a heuristic analysis based on three metrics: frequency (F) measures the total number of times a variable appears; complexity (C) counts involvement in high-cost nonlinear operations (\eg, powers, products of variables, modulo); and coupling (P) represents the number of distinct constraints in which the variable participates. A higher total ``Impact Score'' ($I = F+C+P$) indicates a greater potential for simplification. Applying this analysis to our example, variable $\mathbf{z}$ emerges as the most influential with the highest Impact Score of 11 \wyc{should give the formula, \eg, 4+3+4=11.}, as it is the only variable present in all four constraints and is involved in the most complex operations. This makes it the prime candidate for concretization. Detailed per-variable scores are provided in Appendix~\ref{sec:appendix-motivating-scores}.

\wyc{to discuss: the static heuristic is very important, we should emphasize it and clearly explain why the heuristic choose z=1.}
Once the target variable $\mathbf{z}$ is identified, the next challenge is choosing a value for concretization. From a human expert's perspective, a deeper, structural analysis would be performed. They would recognize that $\mathbf{z}$'s primary contribution to intractability comes from its role within multiple, nested modulo expressions in constraints C3 and C4, which introduce severe discontinuities into the search space. The strategic goal is therefore to select a value that resolves these expressions. The value $1$ is an excellent heuristic choice, as it acts as a neutral element in multiplication and exponentiation. This choice surgically simplifies the formula: it collapses the complex modulo terms (\eg, $(z+15)(z-30)(z-11) \pmod 9$ becomes a constant $5$), while preserving crucial algebraic structure in other constraints by reducing terms like $z^n$ to $1$ and $z \times k$ to $k$, rather than nullifying entire clauses as $z=0$ might. With the intractable complexity originating from $z$ eliminated, the formula becomes substantially simpler for Z3, which then finds a satisfying model in under one second on our workstation \wyc{on the same machine}.

\subsection{Limitations of Static Heuristics}

While this example demonstrates the potential of strategic variable concretization, it also reveals the fundamental limitations of static, heuristic-based approaches like the ``Impact Score''. Such heuristics are often brittle because they rely on superficial proxies. Although our heuristic correctly identified $\mathbf{z}$ as critical, it lacks the capacity for the deeper, structural reasoning required to select an effective value. It cannot understand, as a human expert can, that concretizing $\mathbf{z}$ to $1$ is powerful precisely \wyc{powerful precisely?} because it resolves the modulo-induced discontinuities while preserving algebraic structure elsewhere. This reasoning gap is the core challenge: static heuristics cannot devise a simplification strategy. Static scores provide only weak proxy signals; they cannot perform the semantics-aware value reasoning needed for impactful concretization.

These limitations highlight that a ``one-size-fits-all'' heuristic is insufficient. To overcome these challenges, our work pioneers a framework that learns to reason about solver strategy. We leverage RL to develop a robust policy for variable selection, moving beyond static scores to identify strategic leverage points. To address the challenge of value selection, we utilize LLMs to generate semantics-guided concrete values, replacing the need for the manual, expert-driven heuristic analysis demonstrated above. Crucially, these components operate within a closed-loop system where the agent continuously refines its strategy based on direct feedback from the solver, allowing it to adapt dynamically to the evolving state of the constraint-solving process.

% \paragraph{Takeaways and Roadmap.} This example motivates our learning-based formulation: Section~\ref{sec:problem-formalization} formalizes simplification as an \MDP{}, and Section~\ref{sec:methodology} details how \tool operationalizes the expert-like strategy via an \RL{}+\LLM{} co-pilot with solver-in-the-loop verification.
